\documentclass[10pt]{article}
\usepackage{icmlko}

\icmlmeta{IP 파운데이션 월드모델}{265}{천장을 넘었다}{여지가 처음으로 음수다 --- 그리고 그것은 눈금이 다 됐다는 뜻이기도 하다}

\begin{document}

\twocolumn[
\icmlheader

\begin{abstract}
노트 262 · 263이 새는 축 둘을 막았으니 천장을 다시 재야 했다. 노트 252가
만든 눈금 그대로 --- 유보 \emph{안}에서 5겹 CV, 섞음 귀무로 과적합 바닥 확인.
\textbf{여지가 처음으로 음수다}: 판 0.4842 대 천장 0.4762, \textbf{$-$0.0080}.
도메인별로는 \textbf{아홉 중 여섯이 음수}이고(게임 $-$0.035 · 애니 $-$0.038 ·
세계애니 $-$0.035 · 펀딩 $-$0.035 · 모바일 $-$0.013 · 웹툰 $-$0.009), 양수인
셋은 전부 표본이 작거나 섞음 귀무가 높아 못 읽는 것들이다(도서 163건 ·
팝업 59건 섞음 0.53 · 아이돌 25건 섞음 0.66). \textbf{웹툰이 특히 크게
뒤집혔다} --- 21축에서 $+$0.0552 였는데 지금 $-$0.0091 이다. 두 걸음이 겹친
결과다: 노트 260의 메타 축이 웹툰을 올렸고, 노트 262 · 263의 막기가 웹툰의
\emph{천장}을 0.4495에서 0.4327로 내렸다(막힌 두 축은 CV 모형도 못 쓴다).
\textbf{읽는 법을 조심해야 한다.} 이 천장은 ``도달 가능한 최선''이 아니라
``평범한 능형이 유보 라벨을 보고 도메인마다 따로 배우면 얼마''다. 판은
능형과 나무를 섞고 도메인별 풀링까지 고른 모형이므로, 넘어선 것이 놀랍지는
않다. 정직한 결론은 둘이다 --- \textbf{(가) 풀링이 도메인 전용보다 낫다는
노트 252의 관측이 더 강해졌고}, \textbf{(나) 이 눈금은 이제 어디를 파야
하는지 못 알려 준다.} 노트 252부터 세 번 쓴 ``여지 $+$0.012'' 나침반이
다 됐다.
\end{abstract}
]

\section{다시 재야 했다}

노트 262가 웹툰 \texttt{entry\_friction}(\texttt{daily\_pass})을, 노트 263이
\texttt{goods\_scale}(\texttt{n\_episode})을 막았다. 천장은 축 집합에 딸린
수이므로 다시 재야 한다.

\begin{center}\small
\begin{tabular}{lrrr}
\hline
 & 판 & 천장(CV) & 여지 \\
\hline
17축 (노트 252) & 0.4569 & 0.4697 & $+$0.0128 \\
21축 (노트 260) & 0.4688 & 0.4808 & $+$0.0121 \\
\textbf{23축 (지금)} & \textbf{0.4842} & \textbf{0.4762} &
  $\mathbf{-0.0080}$ \\
\hline
\end{tabular}
\end{center}

\textbf{처음으로 음수다.} 노트 252 · 260에서 두 번 다 $+$0.012 로 붙어 있던
것이 뒤집혔다.

\section{여섯이 제 천장 위에 있다}

\begin{figure}[t]
\centering
\includegraphics[width=\columnwidth]{figs/ceiling.pdf}
\caption{왼쪽 --- 도메인마다 지금(점)과 유보 안 CV(막대). 오른쪽 --- 세 걸음.}
\end{figure}

\begin{center}\small
\begin{tabular}{lrrrr}
\hline
도메인 & $n$ & 지금 & CV & 여지 \\
\hline
애니 & 606 & 0.5365 & 0.4990 & $-$0.0375 \\
게임 & 180 & 0.6355 & 0.6005 & $-$0.0350 \\
펀딩 & 80 & 0.2252 & 0.1903 & $-$0.0349 \\
세계애니 & 300 & 0.5214 & 0.4869 & $-$0.0345 \\
모바일 & 441 & 0.5248 & 0.5118 & $-$0.0130 \\
\textbf{웹툰} & 711 & 0.4418 & 0.4327 & $\mathbf{-0.0091}$ \\
\hline
팝업 & 59 & 0.3346 & 0.4231 & $+$0.0886 \\
도서 & 163 & 0.3496 & 0.4980 & $+$0.1484 \\
아이돌 & 25 & 0.2331 & 0.4064 & $+$0.1733 \\
\hline
\end{tabular}
\end{center}

음수인 여섯이 표본 큰 쪽이고, 양수인 셋은 \textbf{전부 못 읽는 것들}이다 ---
아이돌은 섞은 라벨로도 0.658, 팝업은 0.533 이 나온다(노트 248의 규약대로
읽으면 그 CV 값 자체가 과적합이다). 도서는 노트 253이 이미 닫았다(제 학습이
해로운 유일한 도메인).

\textbf{웹툰이 제일 크게 뒤집혔다} --- 21축에서 $+$0.0552 였는데 $-$0.0091 이다.
두 걸음이 겹쳤다: 노트 260의 메타 축이 웹툰 $\rho$ 를 올렸고(0.3943 $\to$
0.4418), 노트 262 · 263의 막기가 웹툰의 \emph{천장}을 내렸다(0.4495 $\to$
0.4327 --- 막힌 두 축은 CV 모형도 못 쓴다). 위에서 올라오고 아래에서
내려갔다.

\section{이 수를 어떻게 읽어야 하나}

``판이 천장을 넘었다''는 말은 그대로 읽으면 안 된다. 이 천장의 정의는
\textbf{도달 가능한 최선}이 아니다.

\begin{quote}\small
\textbf{평범한 능형}(alpha$=$1)이 \textbf{유보 라벨을 보고} 도메인마다
\textbf{따로} 배워 5겹 CV 로 낸 값.
\end{quote}

판은 능형과 배깅 나무를 순위로 섞고(F23) 도메인별 풀링 선택까지 고른
모형이다. 더 좋은 학습기가 더 나쁜 학습기를 넘는 것은 놀랄 일이 아니다.
\textbf{천장이라는 이름이 과했다.}

그래도 두 가지는 정직하게 남는다.

\paragraph{(가) 풀링이 이겼다} 노트 252가 ``아홉 중 다섯에서 풀링이 도메인
전용을 이긴다''고 적었는데 지금은 여섯이고, 표본 큰 도메인이 전부 거기
있다. \textbf{유보 라벨을 직접 보고 배운 도메인 전용 모형보다 풀링이 낫다}는
관측이 더 강해졌다. 노트 126의 ``도메인 고유 유연성은 비싸다''가 세 번째
형태로 확인된다.

\paragraph{(나) 나침반이 다 됐다} 이 눈금은 노트 252부터 세 번 ``어디를
파야 하나''를 알려 줬다 --- 웹툰 $+$0.055 를 짚었고 노트 260이 그것을
가져왔다. 이제 음수라 \textbf{아무 데도 안 가리킨다.} 더 센 학습기로 천장을
다시 정의하면 되겠지만, 그러면 그 천장은 ``내가 이미 가진 모형''이 되어
동어반복이다.

\section{그러면 무엇으로 고르나}

지금 남은 나침반은 셋이다.

\begin{center}\small
\begin{tabular}{ll}
\hline
나침반 & 상태 \\
\hline
무리 천장(노트 220) & 산다 --- 빼 봤을 때 증분 \\
도메인 천장(노트 222) & 산다 --- 몫 $\times$ (1$-\rho$) \\
유보 안 CV 여지(252) & \textbf{다 됐다} \\
안 쓴 자료 훑기(255) & 산다 --- 노트 260이 두 번째로 썼다 \\
\hline
\end{tabular}
\end{center}

노트 222의 몫 $\times$ (1$-\rho$)로 다시 보면 웹툰이 여전히 1위다
(0.273 $\times$ 0.558 $=$ 0.152). 다만 그것은 ``얼마나 못 맞히나''이지
``얼마나 더 맞힐 수 있나''가 아니다.

\paragraph{남은 것}
\begin{center}\small
\begin{tabular}{ll}
\hline
갈래 & 왜 \\
\hline
안 쓴 자료 훑기 & 유일하게 살아 있는 ``어디를 파나'' \\
\texttt{SOURCE} 채우기 & 20/45 쌍. 가드 덮음이 그만큼이다 \\
라벨 자체 & 노트 209의 누적 라벨 --- 축만 고쳤다 \\
더 센 천장 & 정의하면 동어반복이 된다. 다른 형태가 있나 \\
\hline
\end{tabular}
\end{center}

\paragraph{규약} 진단 눈금은 \textbf{다 됐을 때 다 됐다고 적는다.} 노트 252의
여지는 세 번 값을 했고 이제 음수다. 음수를 ``더 잘하고 있다''로 읽고 계속
쓰면 그때부터는 나침반이 아니라 장식이다.

\end{document}
